\chapter{Conclusiones}
\label{anexo3:cap:conclusiones}
En este capítulo se presentan las conclusiones de la planificación realizada durante el desarrollo
del proyecto, así como una reflexión sobre la experiencia adquirida y las lecciones aprendidas.

\section{Estadísticas del proyecto}
\label{anexo3:sec:estadisticas}
En total se han completado \textbf{619} puntos de historia distribuidos en \textbf{202} tareas
y \textbf{22} sprints. A excepción de los últimos sprints, en los que se dedicó todo el tiempo
disponible al trabajo de fin de grado, se puede estimar que se ha trabajado una media
jornada diaria de \textbf{4} horas. También se podría estimar que el tiempo trabajado en los cuatro
últimos sprints ha sido de \textbf{8} horas diarias, ya que se dedicó todo el tiempo disponible
al proyecto. Por tanto, el tiempo de desarrollo del proyecto habrían sido unas \textbf{520} horas
diarias. Esto implicaría que cada punto de historia completado ha requerido aproximadamente
\textbf{50} minutos de trabajo.

Esta estimación, aunque no es del todo precisa, puede servir para tener una idea del esfuerzo
requerido para completar el proyecto. Se observa que el esfuerzo representado por un punto de historia
se ha ido reduciendo a medida que se ha ido avanzando en el proyecto (algo a lo que he decidido
denominar \enquote{inflación de puntos de historia}). Esto se debe a varias razones, siendo la principal
la falta de experiencia en el desarrollo y planificación de proyectos de software. También ha
podido influir el hecho de que conforme se fue acercando el final del proyecto, el tamaño de las
tareas se fue incrementando, lo que influenció en los errores cometidos, como se verá en el
siguiente apartado.

\section{Lecciones aprendidas}
\label{anexo3:sec:lecciones-aprendidas}
En esta sección se presentan las lecciones aprendidas durante el desarrollo del proyecto en relación
con la planificación y la metodología utilizada. Estas lecciones se centran en los problemas
encontrados durante el desarrollo del sistema, específicamente con la estimación de puntos de
historia y el flujo de trabajo en Scrum. Por tanto, se mencionarán mejores prácticas a la hora
de estimar el esfuerzo necesario para completar las tareas y se reflexionará sobre la metodología
escogida, así como las posibles alternativas que podrían haber sido más adecuadas para un proyecto
de desarrollo individual.

\subsection{Puntos de historia}
Un problema frecuente durante la planificación del proyecto ha sido la estimación de los puntos
de historia. En varias ocasiones se ha dudado sobre el coste que debería tener una tarea. Aunque
es sencillo determinar si este debería ser uno, dos o tres puntos, es más complicado determinar
si una tarea debería tener siete u ocho. Al planificar tareas más grandes, se dudaba si el valor
escogido era el correcto o si debería ser ligeramente mayor o menor. Esto se debe a la Ley de
Weber, que establece que la capacidad de percibir diferencias entre dos estímulos es proporcional
a la relación entre la magnitud de los estímulos más que a la magnitud de los estímulos mismos
\parencite{WeberLaw}.

Esto ha provocado, entre otras cosas, que aunque el trabajo realizado en dos sprints sea el mismo,
el número de puntos de historia completados sea ligeramente diferentes. Una posible solución para
este problema sería utilizar una escala más reducida, como la de Fibonacci, que declara solo
costes específicos (1, 2, 3, 5, 8, 13, …). En la escala de Fibonacci, los números crecen
aproximadamente un 60\% respecto al anterior, lo que hace que las diferencias entre los números
sean más notables. Esto reduce la ambigüedad en la estimación de los puntos de historia y
permite una mejor planificación del proyecto \parencites{Cohn2025Jun}{MariaMedicion}

Otra manera de abordar este problema es utilizar un sistema de medición que no estuviese basado
en puntos de historia, sino en horas aproximadas de trabajo. Uno de estos sistemas es el denominado
\textit{T-shirt sizing}, que utiliza tallas de camisetas (XS, S, M, L, XL) para estimar el tamaño
de las tareas. Este sistema es rápido y fácil de usar, ya que no requiere una escala numérica
específica. Elimina el problema de la ambigüedad en la estimación de los puntos de historia, puesto
que solo existen cinco tallas diferentes. Sin embargo, puede ser menos preciso que otros sistemas,
ya que no proporciona una medida cuantitativa del esfuerzo requerido para completar una tarea y
carecen de la resolución necesaria para tareas más complejas. Por esto, es recomendable mantener
la cantidad de tiempo necesaria por tarea reducida y dividir las tareas más grandes en subtareas
más pequeñas y manejables \parencites{mallidi2021study}{Krdzic2025May}.

\subsection{Metodología usada}
Este ha sido el primer proyecto que he desarrollado utilizando una metodología ágil.
La elección de este enfoque y no de uno predictivo ha permitido una mayor flexibilidad y adaptabilidad
al cambio, lo que ha sido especialmente útil en un proyecto de desarrollo como este, donde es
común encontrar problemas técnicos inesperados que afectarían enormemente a una planificación
predictiva. Como ya se mencionó en el capítulo que describe el modelo de planificación escogido
(ver \nameref{anexo3:cap:modelo}), planificar de manera predictiva un proyecto de gran tamaño
como es este es un proceso complejo para el cual no se posee la experiencia necesaria.

Sin embargo, la elección de Scrum sobre otras metodologías ágiles ha supuesto algunos retos. Por
ejemplo, la planificación de sprints ha requerido una buena cantidad de tiempo y esfuerzo,
especialmente al principio del proyecto, cuando se estaba aprendiendo a utilizar la herramienta
de planificación. Además, puesto que una vez iniciado un sprint no se puede modificar el conjunto
de tareas a realizar, ocasionalmente se descubre un error o un problema técnico que requiere una
modificación en la planificación.

Otra característica de Scrum que es importante mencionar es la necesidad de realizar reuniones
diarias y la existencia de roles específicos dentro del equipo. Aunque estas dos características
son útiles para mantener al equipo alineado y enfocado en los objetivos del proyecto, son algo
que no se ha podido aplicar en este proyecto, ya que se ha desarrollado de manera individual.

Por tanto, aunque Scrum es una metodología ágil muy útil para proyectos de desarrollo, su
aplicación en un proyecto individual puede no ser la más adecuada. Una alternativa podría
ser utilizar una metodología ágil más flexible, como Kanban, que no requiere la planificación
de sprints ni la existencia de roles específicos dentro del equipo. Kanban permite una mayor
adaptabilidad al cambio y una mayor flexibilidad en la planificación, lo que puede ser más adecuado
para proyectos individuales \parencite{junior2010variations}.

\diagrama{
    nombre = {Pizarra de Kanban},
    fuente = {Chris Huffman, CC en \url{https://flic.kr/p/4yvFP2}},
    etiqueta = {anexo3:fig:kanbanconclusiones},
    archivo = {anexo3/imgs/otrokanban.png},
}

Puede que Kanban sea una opción demasiado laxa para un proyecto de desarrollo de gran tamaño en el
que existen fechas de entrega y objetivos específicos. Para solucionar estos problemas se podría
haber utilizado una metodología ágil híbrida, que combine las ventajas de Scrum y Kanban. Por ejemplo,
Scrumban es una metodología que combina la planificación de sprints de Scrum con la flexibilidad
de Kanban. Scrumban carece de roles específicos y reduce el número de reuniones diarias, así como
añadir a la planificación el flujo continuo de trabajo propio de Kanban \parencites{Stoica2016AnalyzingAD}%
{Scrumban}. Es por esto que se cree que esta metodología, que combina las ventajas de Scrum y Kanban,
podría haber sido  una mejor opción para este proyecto.

